\documentclass[10pt,a4paper]{amsart}
\usepackage[margin=19mm]{geometry}
\usepackage{amsmath,amssymb,mathtools}
\usepackage[scr=rsfs,cal=euler]{mathalfa}
\usepackage{xfrac}
\usepackage[unicode,hidelinks,bookmarks=false]{hyperref}
\newcommand{\ie}{i.e.\ }
\newcommand{\cf}{cf.\ }
\newcommand{\resp}{respectively\ }
\newcommand{\Subsection}{\subsection}
\providecommand{\nobreakdash}{\nobreak\hbox{-}}
\setlength{\emergencystretch}{2em}
\multlinegap0pt
\usepackage{bm}
\usepackage{bbm}
\usepackage{dsfont}
\usepackage{xspace}
\usepackage{cancel}
\usepackage{mathtools}
\usepackage{amsthm}
\makeatletter
\@ifundefined{thm*}{\newtheorem*{thm*}{Theorem}}{}
\makeatother
\DeclareMathOperator{\clh}{\overline{span}}
\DeclareMathOperator{\lh}{span}
\newcommand{\id}{\mathrm{id}}
\title[Total sequences and the lower frame inequality]{Total sequences and the lower frame inequality}
\author{Narutaka Ozawa}
\date{Source: arXiv:2507.04899v2 (CC BY 4.0)}
\begin{document}
\maketitle

\noindent\emph{Source and licence.} Total sequences and the lower frame inequality. Version arXiv:2507.04899v2 (CC BY 4.0), distributed under the Creative Commons Attribution 4.0 International licence, as recorded in the arXiv licence metadata for this exact version and shown on the abstract page https://arxiv.org/abs/2507.04899v2. Full rights notice: licenses/arxiv-2507.04899-CC-BY-4.0.txt.\par\smallskip

\noindent\textit{Editorial scope.} Counted unit: the Thm* whose source label is \texttt{None} in the article (source0.tex lines 72--86, source label \texttt{None}), delivered together with the article's own complete original proof of it (source0.tex lines 88--136, one \texttt{proof} environment of 49 lines). No mathematics is written, repaired or re-derived: statement and proof are the article's own text line for line. Required context retained uncounted: none. Every definition, display and label of those lines is retained unchanged. Omitted from this package and not redistributed: 1--71, 87--87, 137--248. Nothing inside a retained excerpt is omitted or altered: every retained body is delivered exactly as the article's lines are written, so no omission receipt of any kind accompanies this package. Citations: the retained excerpts carry no citation command at all, so no bibliography entry of the article is transcribed and no reference is redistributed. Reference adaptation: none was needed and none was made; every reference inside the delivered unit resolves inside it, so no unresolved reference or citation is produced. Printed slips kept literal: no printed word, symbol, hypothesis or number of the counted statement or of its proof was altered. Layout and class: the article's own class is \texttt{amsart} with packages including hyperref, mathtools; the wrapper is the lane's pinned \texttt{amsart} editorial wrapper at 10pt on A4, which restates the theorem-like environments the retained text needs and adds the article's own macro definitions that the retained text needs (\texttt{\textbackslash clh}, \texttt{\textbackslash id}, \texttt{\textbackslash lh}), with extra wrapper packages (bm, bbm, dsfont, xspace, cancel, mathtools). The delivered numbering is the unit's own. Assets: no retained span contains a figure environment, a graphics-inclusion command, a PDF-page inclusion, a table or a TikZ picture; no image file is copied into this package, the package declares no asset dependency and no image file is redistributed with it. Licence: version arXiv:2507.04899v2 of this article is distributed under the Creative Commons Attribution 4.0 International licence, as recorded in the arXiv licence metadata for this exact version and shown on the abstract page https://arxiv.org/abs/2507.04899v2. A full rights notice is delivered at licenses/arxiv-2507.04899-CC-BY-4.0.txt. Characters: every retained excerpt is pure ASCII, so the delivered engine, XeLaTeX with its default Latin Modern text font, reports no missing character.
\begin{thm*}
Let $( v_n )_{n=1}^\infty$ be a total sequence in a Hilbert space $H$. 
Then, there are sequences $(w_n)_{n=1}^\infty$ and $(z_n)_{n=1}^\infty$ 
in $H$ such that $z_n\in\lh\{ v_k : k\geq n/2\}$ and 
$$\id_H=\sum_{n=1}^\infty w_nz_n^*$$
in the strong operator topology. 
In particular, 
there is a sequence $(\lambda_n)_{n=1}^\infty$ 
of positive numbers such that $(\lambda_n^{1/2}v_n)_{n=1}^\infty$
satisfies the lower frame inequality: 
\[
\| x\|^2 \le \sum_{n=1}^\infty\lambda_n |v_n^*x|^2 
\]
for every $x\in H$, with the right hand side being possibly $\infty$.
\end{thm*}

\begin{proof}
Let $H_n\coloneqq\clh\{ v_k : k\geq n \}$ and $H_\infty\coloneqq\bigcap_n H_n$. 
Then $H=H_1\supset H_2\supset\cdots$ and 
$\dim(H_n\cap H_{n+1}^\perp)\le1$. 
For each $n$, if $H_n=H_{n+1}$, then set $x_n\coloneqq0$; 
else take a unit vector 
$x_n\in H_n\cap H_{n+1}^\perp$. 
Thus $\{ x_n : x_n\neq0\}$ is an orthonormal 
basis for $H_\infty^\perp$.
We also fix an orthonormal basis $\{ y_n\}$ for $H_\infty$, 
which is possibly empty, and 
set $u_{2n-1} \coloneqq x_n$ and $u_{2n}\coloneqq y_n$ 
or 0 if $y_n$ is not defined. 
Then $\{ u_n : u_n\neq0\}$ is an orthonormal basis 
for $H$ such that $u_n\in H_{\lceil n/2\rceil}$. 
We choose $z_n\in\lh\{ v_k : k\geq n/2\}$ such that 
$\| u_n - z_n \| < 3^{-n}$. 
Then $T\coloneqq\sum_n u_n (u_n-z_n)^*$ converges absolutely, 
$\|T\|<1$, and 
$$\id_H-T=\bigl(\sum_n u_nu_n^*\bigr)-T = \sum_n u_nz_n^*$$ 
in the strong operator topology. 
Thus $w_n\coloneqq(\id_H-T)^{-1}u_n$ meet our demand. 
This proves the first half of Theorem.

For each $n$, 
write $z_n = \sum_{k\geq n/2} \gamma^{(n)}_k v_k$ 
as a finite sum. 
Note that by the Cauchy--Schwarz inequality, 
$(\sum_{n=1}^\infty \beta_n)^2\le \sum_{n=1}^\infty 2^n\beta_n^2$ 
holds for every sequence $(\beta_n)_{n=1}^\infty$ 
of non-negative numbers, where 
both sides can be $\infty$. 
Thus for every $x\in H$ 
\[
\|x\|^2 = \|\sum_n w_n z_n^*x\|^2 
\le\sum_n 2^n\|w_n\|^2|z_n^*x|^2,
\]
while 
\[
|z_n^*x|^2=|\sum_{k\geq n/2}\gamma^{(n)}_k v_k^*x |^2
\le 2^{-(n/2)+1}\sum_{k\geq n/2} 2^k|\gamma^{(n)}_k|^2 |v_k^*x |^2;
\]
which amounts to 
\[
\|x\|^2 \le \sum_n\sum_{k\geq n/2}2^{(n/2)+k+1}\|w_n\|^2|\gamma^{(n)}_k|^2 |v_k^*x |^2
=\sum_{k=1}^\infty\lambda_k |v_k^*x |^2,
\]
where $\lambda_k\coloneqq\sum_{n=1}^{2k}2^{(n/2)+k+1}\|w_n\|^2|\gamma^{(n)}_k|^2$.
\end{proof}

\end{document}
